% ============================================================
\chapter{Submanifolds and the Second Fundamental Form}
\label{ch:submanifolds}
% ============================================================

\section{Riemannian submanifolds}

\begin{definition}[Riemannian submanifold]
\index{submanifold}
Let $(M, g)$ be a Riemannian manifold of dimension $n$ and $\Sigma \subset M$
a submanifold of dimension $k < n$. The \emph{induced metric} on $\Sigma$ is
$\bar{g} = \iota^* g$, where $\iota : \Sigma \hookrightarrow M$ is the inclusion.
The pair $(\Sigma, \bar{g})$ is called a \emph{Riemannian submanifold}.
\end{definition}

\begin{definition}[Tangent bundle decomposition]
\index{normal bundle}
At each point $p \in \Sigma$, the tangent bundle of $M$ decomposes orthogonally:
\[
    T_pM = T_p\Sigma \oplus N_p\Sigma,
\]
where $N_p\Sigma = (T_p\Sigma)^\perp$ is the \emph{normal bundle} of $\Sigma$ in $M$.
We write $(\cdot)^\top$ and $(\cdot)^\perp$ for the corresponding projections.
\end{definition}

\begin{intuition}
Think of a surface $\Sigma^2$ in $\R^3$. At each point, the tangent space of $\R^3$
decomposes into the tangent plane of $\Sigma$ and the normal direction $\nu$. The
second fundamental form measures how $\Sigma$ ``bends'' inside $\R^3$, i.e., how
the normal $\nu$ varies along $\Sigma$.
\end{intuition}

\section{Induced connection and second fundamental form}

\begin{definition}[Gauss formula]
\index{Gauss formula}
Let $\nabla$ be the Levi-Civita connection of $(M, g)$ and $\bar{\nabla}$ that of
$(\Sigma, \bar{g})$. For all vector fields $X, Y$ tangent to $\Sigma$:
\[
    \nabla_X Y = \bar{\nabla}_X Y + \mathrm{I\!I}(X, Y),
\]
where $\bar{\nabla}_X Y = (\nabla_X Y)^\top$ and $\mathrm{I\!I}(X, Y) = (\nabla_X Y)^\perp$.
\end{definition}

\begin{definition}[Second fundamental form]
\index{second fundamental form}
The \emph{second fundamental form} is the symmetric bilinear map:
\[
    \mathrm{I\!I} : \mathfrak{X}(\Sigma) \times \mathfrak{X}(\Sigma) \to \Gamma(N\Sigma),
    \quad \mathrm{I\!I}(X, Y) = (\nabla_X Y)^\perp.
\]
\end{definition}

\begin{proposition}[Symmetry of $\mathrm{I\!I}$]
The second fundamental form is symmetric: $\mathrm{I\!I}(X, Y) = \mathrm{I\!I}(Y, X)$.
\end{proposition}

\begin{proof}
We have $\mathrm{I\!I}(X, Y) - \mathrm{I\!I}(Y, X) = (\nabla_X Y - \nabla_Y X)^\perp
= [X, Y]^\perp$. Since $X$ and $Y$ are tangent to $\Sigma$, the Lie bracket
$[X, Y]$ is also tangent, so $[X, Y]^\perp = 0$.
\end{proof}

\begin{definition}[Shape operator (Weingarten map)]
\index{shape operator}
For a normal field $\xi \in \Gamma(N\Sigma)$, the \emph{shape operator} is the
endomorphism $A_\xi : T\Sigma \to T\Sigma$ defined by:
\[
    A_\xi(X) = -(\nabla_X \xi)^\top.
\]
It is related to $\mathrm{I\!I}$ by:
$g(\mathrm{I\!I}(X, Y), \xi) = g(A_\xi(X), Y)$.
\end{definition}

\begin{proposition}[Self-adjointness of $A_\xi$]
The shape operator is self-adjoint:
$\bar{g}(A_\xi(X), Y) = \bar{g}(X, A_\xi(Y))$.
\end{proposition}

\begin{proof}
$\bar{g}(A_\xi(X), Y) = g(\mathrm{I\!I}(X, Y), \xi) = g(\mathrm{I\!I}(Y, X), \xi)
= \bar{g}(A_\xi(Y), X)$, by symmetry of $\mathrm{I\!I}$.
\end{proof}

\section{Mean curvature and minimal submanifolds}

\begin{definition}[Mean curvature]
\index{mean curvature}
Let $(e_1, \ldots, e_k)$ be an orthonormal basis of $T_p\Sigma$. The \emph{mean
curvature vector} is:
\[
    \vec{H} = \frac{1}{k} \sum_{i=1}^{k} \mathrm{I\!I}(e_i, e_i)
    = \frac{1}{k} \tr(\mathrm{I\!I}).
\]
\end{definition}

\begin{definition}[Minimal submanifold]
\index{minimal submanifold}
A submanifold $\Sigma$ is \emph{minimal} if $\vec{H} = 0$, i.e., $\tr(\mathrm{I\!I}) = 0$.
\end{definition}

\begin{remark}
The term ``minimal'' is historical. Minimal submanifolds are critical points of
the volume functional: $\delta \operatorname{Vol}(\Sigma) = 0$. They do not
necessarily minimize volume.
\end{remark}

\begin{example}[Catenoid]
\index{catenoid}
The surface of revolution in $\R^3$ parametrized by:
\[
    \Phi(u, v) = (\cosh u \cos v, \cosh u \sin v, u), \quad (u, v) \in \R \times [0, 2\pi),
\]
is a minimal surface with $H = 0$.
\end{example}

\begin{example}[Sphere $S^{n-1} \subset \R^n$]
The sphere $S^{n-1}(r)$ of radius $r$ has second fundamental form
$\mathrm{I\!I}(X, Y) = -\frac{1}{r} g(X, Y)\, \nu$, where $\nu$ is the outward
unit normal. The mean curvature is $\vec{H} = -\frac{1}{r}\nu$: the sphere is
not minimal.
\end{example}

\section{Gauss and Codazzi equations}

\begin{theorem}[Gauss equation]
\label{thm:gauss}
\index{Gauss equation}
Let $X, Y, Z, W$ be vector fields tangent to $\Sigma$. The curvature of $\Sigma$
is related to that of $M$ by:
\begin{multline*}
    \bar{g}(\bar{R}(X, Y)Z, W) = g(R(X, Y)Z, W) \\
    + g(\mathrm{I\!I}(X, Z), \mathrm{I\!I}(Y, W))
    - g(\mathrm{I\!I}(Y, Z), \mathrm{I\!I}(X, W)).
\end{multline*}
\end{theorem}

\begin{proof}
By definition, $\bar{R}(X,Y)Z = \bar{\nabla}_X \bar{\nabla}_Y Z
- \bar{\nabla}_Y \bar{\nabla}_X Z - \bar{\nabla}_{[X,Y]} Z$. Using the Gauss
formula $\nabla_X Y = \bar{\nabla}_X Y + \mathrm{I\!I}(X, Y)$ and expanding:
\begin{align*}
    \nabla_X(\bar{\nabla}_Y Z + \mathrm{I\!I}(Y,Z))
    &= \bar{\nabla}_X \bar{\nabla}_Y Z + \mathrm{I\!I}(X, \bar{\nabla}_Y Z) \\
    &\quad - A_{\mathrm{I\!I}(Y,Z)} X + \nabla_X^\perp \mathrm{I\!I}(Y,Z).
\end{align*}
Antisymmetrizing in $X, Y$ and projecting onto the tangential component yields:
\[
    (\bar{R}(X,Y)Z)^\top = (R(X,Y)Z)^\top + A_{\mathrm{I\!I}(X,Z)}Y - A_{\mathrm{I\!I}(Y,Z)}X.
\]
Taking the inner product with $W \in T\Sigma$ and using
$g(A_\xi Y, W) = g(\mathrm{I\!I}(Y,W), \xi)$ gives the result.
\end{proof}

\begin{corollary}[Gauss's Theorema Egregium]
\index{Theorema Egregium}
For a surface $\Sigma^2 \subset \R^3$, the Gaussian curvature $K$ is intrinsic:
\[
    K = \kappa_1 \kappa_2 = \det(A_\nu),
\]
where $\kappa_1, \kappa_2$ are the principal curvatures.
\end{corollary}

\begin{theorem}[Codazzi equation]
\label{thm:codazzi}
\index{Codazzi equation}
For all fields $X, Y, Z$ tangent to $\Sigma$ and $\xi$ normal:
\[
    g(R(X, Y)Z, \xi) = g\bigl((\bar{\nabla}_X \mathrm{I\!I})(Y, Z)
    - (\bar{\nabla}_Y \mathrm{I\!I})(X, Z),\, \xi\bigr),
\]
where $(\bar{\nabla}_X \mathrm{I\!I})(Y, Z) = \nabla_X^\perp(\mathrm{I\!I}(Y, Z))
- \mathrm{I\!I}(\bar{\nabla}_X Y, Z) - \mathrm{I\!I}(Y, \bar{\nabla}_X Z)$.
\end{theorem}

\begin{remark}
When $M = \R^n$ (zero curvature), the Codazzi equation simplifies to:
$(\bar{\nabla}_X \mathrm{I\!I})(Y,Z) = (\bar{\nabla}_Y \mathrm{I\!I})(X,Z)$.
\end{remark}

\section{Hypersurfaces}

\begin{definition}[Hypersurface]
\index{hypersurface}
A \emph{hypersurface} is a submanifold of codimension $1$. The normal bundle has
rank $1$, generated by a unit field $\nu$.
\end{definition}

For a hypersurface, the scalar second fundamental form is:
\[
    h(X, Y) = g(\mathrm{I\!I}(X, Y), \nu) = g(A_\nu(X), Y).
\]

\begin{definition}[Principal curvatures]
\index{principal curvatures}
The \emph{principal curvatures} $\kappa_1, \ldots, \kappa_{n-1}$ are the eigenvalues
of $A_\nu$. The eigenvectors are the \emph{principal directions}.
\end{definition}

\begin{proposition}
For a hypersurface:
\begin{align*}
    H &= \frac{1}{n-1}(\kappa_1 + \cdots + \kappa_{n-1})
       = \frac{1}{n-1} \tr(A_\nu), \\
    K_{\text{Gauss}} &= \kappa_1 \cdots \kappa_{n-1} = \det(A_\nu).
\end{align*}
\end{proposition}

\section{Totally geodesic and umbilical submanifolds}

\begin{definition}[Totally geodesic]
\index{totally geodesic}
$\Sigma$ is \emph{totally geodesic} if $\mathrm{I\!I} \equiv 0$.
\end{definition}

\begin{example}
Great circles of $S^2$ and equatorial spheres $S^k \subset S^n$ are totally geodesic.
\end{example}

\begin{definition}[Totally umbilical]
\index{totally umbilical}
$\Sigma$ is \emph{totally umbilical} if $\mathrm{I\!I}(X, Y) = g(X, Y)\, \vec{H}$.
\end{definition}

\begin{example}
The spheres $S^{n-1}(r) \subset \R^n$ are totally umbilical, with
$\kappa_1 = \cdots = \kappa_{n-1} = 1/r$.
\end{example}

\begin{keyformulas}[title=\textbf{Fundamental formulas for submanifolds}]
\begin{align*}
    \text{Gauss:} \quad & \nabla_X Y = \bar{\nabla}_X Y + \mathrm{I\!I}(X, Y) \\
    \text{Weingarten:} \quad & \nabla_X \xi = -A_\xi(X) + \nabla_X^\perp \xi \\
    \text{Gauss (curv.):} \quad & \bar{g}(\bar{R}(X,Y)Z, W) = g(R(X,Y)Z, W)
        + g(\mathrm{I\!I}(X,Z), \mathrm{I\!I}(Y,W)) \\
    & \qquad\qquad\qquad\quad - g(\mathrm{I\!I}(Y,Z), \mathrm{I\!I}(X,W)) \\
    \text{Mean curv.:} \quad & \vec{H} = \frac{1}{k}\tr(\mathrm{I\!I})
\end{align*}
\end{keyformulas}

\section{Exercises}

\begin{exercise}
Compute the second fundamental form of the torus of revolution
$\Phi(\theta, \phi) = ((R + r\cos\theta)\cos\phi,
(R + r\cos\theta)\sin\phi, r\sin\theta)$ in $\R^3$. Determine the principal
curvatures and mean curvature.
\end{exercise}

\begin{exercise}
Show that for a hypersurface $\Sigma^{n-1} \subset \R^n$, the Gauss equation gives
$\bar{K}(\sigma) = \kappa_i \kappa_j$ for the plane $\sigma$ spanned by principal
directions $e_i$ and $e_j$.
\end{exercise}

\begin{exercise}
Let $\Sigma$ be a minimal surface in $\R^3$. Show that $\kappa_1 = -\kappa_2$ at
every point, and that $K = -\kappa_1^2 \leq 0$.
\end{exercise}

\begin{exercise}
Show that the helicoid $\Phi(u,v) = (v\cos u, v\sin u, au)$ is a minimal surface
in $\R^3$.
\end{exercise}

\begin{exercise}
Let $\Sigma$ be a compact boundaryless hypersurface of $\R^n$. Show that there
exists a point $p \in \Sigma$ where all principal curvatures are strictly positive.
\emph{Hint:} consider the point farthest from the origin.
\end{exercise}

\begin{exercise}
Let $f : \R^2 \to \R$ be a smooth function and $\Sigma = \{(x, y, f(x,y))\}$ the
graph of $f$ in $\R^3$. Compute the second fundamental form, principal curvatures,
and mean curvature of $\Sigma$ in terms of $f$ and its derivatives.
\end{exercise}
